\chapter{SSIM, noise and timing}\label{ch:noise}

In the meeting of 18 November the group asked for four things before the evaluation: report the metrics that Reed et
al.\ and the group's earlier work use (SSIM and RMSE), show what happens as the noise grows, compare our SSIM with
the paper's, and say how fast each method is. This chapter answers all four. Nothing was trained and nothing was
chosen on COMSOL; every number comes from models and settings fixed earlier.

One decision was taken the same day: \textbf{COMSOL is now reported only on the re-solved L1\_C1} (the 1e-8
tolerance solve of chapter~\ref{ch:comsol}). The original exports are known to be wrong, so they appear only as a
reproduction check. Validation uses the free-space simulator.

\section{What SSIM measures}

Relative L2 adds up every pixel's error. SSIM (Wang et al.\ 2004) instead compares two images in small windows and
asks three questions per window: is the mean the same (luminance), is the spread the same (contrast), and do the
two go up and down together (structure)? With local means $\mu_x,\mu_y$, variances $\sigma_x^2,\sigma_y^2$ and
covariance $\sigma_{xy}$, computed with an $11\times11$ Gaussian window of width 1.5 pixels,
\begin{equation}
\mathrm{SSIM} = \frac{(2\mu_x\mu_y + c_1)(2\sigma_{xy} + c_2)}{(\mu_x^2+\mu_y^2+c_1)(\sigma_x^2+\sigma_y^2+c_2)},
\qquad c_1 = (0.01R)^2,\; c_2 = (0.03R)^2 ,
\end{equation}
averaged over all windows. $R$ is the data range, here the target's max $-$ min. SSIM is 1 for identical images. It is
the function \texttt{ssim} in \texttt{scripts/magic\_unet.py}; the same function scores both $B$ and $J$.

Two things to keep in mind. SSIM is dominated by the smooth parts of an image, which every method gets right, so the
differences between methods look small (0.966 against 0.960). And SSIM forgives a wrong overall amplitude more than L2
does: a map with the right shape but 60\% of the amplitude still has high structure.

\section{Every method on the re-solved COMSOL coil}

Each method was run again on the re-solved pair (code of amendment A7, with SSIM and RMSE added). Before scoring,
each method reproduces its saved score on the original L1\_C1; every check passes (worst $4\times10^{-5}$) except
FISTA-Net's, which misses the $10^{-4}$ tolerance by the same $3$--$4\times10^{-4}$ as before.

\begin{center}
\small
\begin{tabular}{lrrrrrr}
\toprule
 & \multicolumn{3}{c}{3\,$\to$\,0.5\um} & \multicolumn{3}{c}{10\,$\to$\,0.5\um} \\
Method & L2 & SSIM & RMSE (\textmu T) & L2 & SSIM & RMSE (\textmu T) \\
\midrule
MAGIC-UNet + correction & \textbf{22.0\%} & \textbf{0.966} & \textbf{4.17} & \textbf{26.2\%} & \textbf{0.965} & \textbf{5.02} \\
MAGIC-UNet raw & 24.6\% & 0.963 & 4.63 & 26.3\% & 0.963 & 5.03 \\
FISTA (925 / 3{,}000 it.) & 25.8\% & 0.960 & 4.97 & 47.2\% & 0.928 & 8.98 \\
CSC-LISTA$^\ast$ & 34.1\% & 0.948 & 6.50 & 98.0\% & 0.868 & 18.5 \\
FISTA-Net$^\ast$ & 32.6\% & 0.930 & 6.26 & 90.6\% & 0.881 & 17.1 \\
Tikhonov & 48.8\% & 0.946 & 9.22 & 98.1\% & 0.871 & 18.5 \\
\bottomrule
\end{tabular}

\smallskip
{\footnotesize $^\ast$ At 10\um with their 3\um settings: not like-for-like (chapter~\ref{ch:ten_results}).
One coil, one seed.}
\end{center}

SSIM ranks the methods exactly as L2 does, at both standoffs. The ranking does not depend on the metric.

\section{The noise sweep}

\paragraph{How it was done.} The first 20 validation layouts of the free-space simulator (the same layouts as every
earlier simulated score), at 10 and 3\um. Each layout keeps its own noise pattern; only its level changes: 0, 1, 2,
5, 10, 25, 50, 67 and 100\% of each component's range. The level ``native'' keeps each layout's own 1--2.5\% and
must reproduce the saved scores; it does, to $1.4\times10^{-5}$ (10\um) and $1.1\times10^{-5}$ (3\um). The
correction strength $\mu$ stays at its value chosen for 1--2.5\% noise.

\paragraph{The paper's noise unit.} Reed et al.\ write the noise as $\sigma_n = \sigma/\max|B|$. On our layouts
1\% of the range is $\sigma_n\approx0.015$, and 67\% is $\sigma_n\approx1$.

\begin{figure}[h]
\centering
\includegraphics[width=\textwidth]{ns_noise_sweep}
\caption{(a) Field error at 0.5\um as the noise grows. (b) SSIM of the current density, the metric the paper reports,
with the paper's two values in grey. Shaded: the noise the networks were trained with.}
\end{figure}

\begin{center}
\small
\begin{tabular}{lrrrrr}
\toprule
Noise & $\sigma_n$ & 10\um: $B$ error & 10\um: $J$ SSIM & 3\um: $B$ error & 3\um: $J$ SSIM \\
\midrule
0\% & 0 & 13.2\% & 0.978 & 5.8\% & 0.982 \\
native 1--2.5\% & 0.024 & 14.7\% & 0.977 & 12.4\% & 0.865 \\
5\% & 0.074 & 21.8\% & 0.957 & 31.2\% & 0.565 \\
10\% & 0.149 & 48.9\% & 0.882 & 49.5\% & 0.332 \\
25\% & 0.37 & 92.7\% & 0.737 & 355\% & 0.119 \\
67\% & 1.0 & 112\% & 0.610 & 859\% & 0.024 \\
\bottomrule
\end{tabular}
\end{center}

\paragraph{What it shows.}
\begin{enumerate}
\item \textbf{Without noise our current SSIM is close to the paper's}: 0.978 and 0.982 against 0.99.
\item \textbf{Beyond the training noise ours falls faster.} At $\sigma_n=1$ the paper reports 0.84; we get 0.61 at
10\um and 0.02 at 3\um. We trained only at $\sigma_n\approx0.015$--$0.04$; the paper does not say what it used. A
network only learns the noise it has seen, so the fix is to train with a wider range of noise.
\item \textbf{The 10\um network is far more robust than the 3\um one.} Run B averages $7\times7$ input pixels into
each 0.7\um network pixel, run 2 only $2\times2$, so the noise it sees is several times smaller.
\item \textbf{At 3\um the noise is already most of the error}: 12.4\% at the native noise against 5.8\% with none.
Better data would help this setting more than a better network.
\item \textbf{The correction must be re-tuned for much noisier data.} Its $\mu$ was chosen at 1--2.5\% noise. At 3\um
with 25\% noise or more it amplifies the noise (355\% against 124\% without it), because $H/(H^2+\mu)$ with a small
$\mu$ boosts every frequency where the data are mostly noise.
\end{enumerate}

\paragraph{Against FISTA at 3\um.} FISTA-925 was run on the first 5 of the layouts at the same noise levels, once with
its $\lambda$ fixed (as MAGIC-UNet's settings are) and once with $\lambda\times1$, 3 and 10, keeping the best per level
(chosen on the same layouts, which favours FISTA).

\begin{center}
\small
\begin{tabular}{lrrr}
\toprule
Noise & MAGIC-UNet + correction & FISTA, $\lambda$ fixed & FISTA, $\lambda$ re-tuned \\
\midrule
0\% & \textbf{6.1\%} & 14.4\% & 14.4\% \\
native 1--2.5\% & \textbf{14.8\%} & 17.8\% & 17.8\% \\
2\% & 16.7\% & 17.9\% & 17.9\% \\
5\% & 30.0\% & 35.6\% & \textbf{27.5\%} \\
10\% & 47.1\% & 88.9\% & \textbf{37.4\%} \\
25\% & 313\% & 350\% & \textbf{92.1\%} \\
\bottomrule
\end{tabular}
\end{center}

Inside the training noise MAGIC-UNet is better or level. Above it the answer depends on who is allowed to adapt: with
everything left at its training-noise setting MAGIC-UNet still wins, but a FISTA whose one parameter $\lambda$ is
raised for the noise wins from 5\% upward. The matching step for the network (re-choosing $\mu$, or training with more
noise) was not done. The lesson is general: a learned prior is only as good as the conditions it was trained for,
while a classical method with a single knob can be re-tuned on the spot.

FISTA at 10\um (5 layouts) was still running when this chapter was written.

\section{How fast}

Wall time per $140\times140$\um image ($1400\times1400$ pixels), one RTX A6000, nothing else running:

\begin{center}
\small
\begin{tabular}{lrr}
\toprule
Method & 3\,$\to$\,0.5\um & 10\,$\to$\,0.5\um \\
\midrule
MAGIC-UNet, network only & 0.42\,s & 0.12\,s \\
MAGIC-UNet + height read-out + correction & 0.79\,s & 0.67\,s \\
FISTA (925 / 3{,}000 iterations) & 91\,s & about 290\,s \\
\bottomrule
\end{tabular}
\end{center}

MAGIC-UNet is about 115 times faster than FISTA at 3\um and 400 times faster at 10\um, and it estimates the standoff, which FISTA has to be told.

\section{The FISTA depth at 10\um}

Every 10\um FISTA number uses 3{,}000 iterations. The validation curve to 10{,}000 iterations has its minimum at
3{,}125, but only 0.004 points lower (65.457\% against 65.460\%). The rule fixed in advance changes the depth only for a
gain of at least 0.5 points, so 3{,}000 stays. In words: the best depth is about 3{,}000, and the two give the same
error.
